\section{A packing bound independent of the positive excess}
\label{sec:four-high-connectivity-excess-free-packing}

The absence of a moderate cut with two large-rank shores bounds the number of small cyclic regions.
The bound below assumes only that the forest probability exceeds one half.
It does not depend on the amount of that excess.

\begin{lemma}[Terminal bipartitions along a forest]
\label{lem:four-terminal-bipartition}
Let $Q$ be a connected loopless multigraph and let $S$ be a nonempty proper vertex set.
For each vertex $v$, let $L_v$ be the span of its incident coordinate forms on $\mathcal Z(Q)$.
Put $U=\sum_{s\in S}L_s$ and $B=|\delta_QS|$.
Among vertices $v\notin S$ with $L_v\cap U\ne0$, at most $2B$ are exceptional.
The others form at most $2B$ ordered classes.
Within each class, the vertices are articulation vertices of one component of $Q-S$ and give the same bipartition of its boundary terminals.
Here each edge to $S$ supplies one distinct terminal.
\end{lemma}

\begin{proof}
Exclude the at most $B$ vertices adjacent to $S$.
In each component of $Q-S$, choose a spanning tree, attach a leaf for each boundary terminal, and retain the union of paths between those leaves.
This forest has at most $B$ leaves and at most $B$ branching vertices.
Exclude its branching vertices and suppress its degree-two vertices.
The resulting forest has at most $2B$ edges.

To justify that these edges contain every other relevant vertex, fix such a vertex $v$ and choose $0\ne a\in L_v\cap U$.
Represent $a$ both by a subset of the star at $v$ and by edges incident with $S$.
Their symmetric difference annihilates the cycle space, so it is a cut.
In the component $K$ of $Q-S$ containing $v$, the star subset therefore joins $v$ to a union of components of $K-v$.
Its functional equals the sum of the terminal coordinates in that union.
The sum of all terminal coordinates of $K$ is zero.
Since $a\ne0$, at least two components of $K-v$ contain terminals.
Every spanning tree must connect those groups through $v$.
After the branching exceptions, exactly two terminal-bearing components remain.
Their terminal sets are the two sides of one edge of the suppressed forest.
Thus vertices on the same edge give the same actual terminal bipartition, as asserted.
\end{proof}

\begin{proposition}
\label{prop:four-excess-free-cycle-packing}
Let $G$ be connected and cubic, and let $P$ be a marked vertex set.
Fix integers $d\ge2$ and $H\ge1$.
Assume that no cut of order at most $2d$ has two connected shores, each of rank at least $H$.
If
\[
 \Pr(\xi_G-P\text{ is a forest})>\frac12,
\]
then every family of $q$ vertex-disjoint connected cyclic regions avoiding $P$, each with boundary at most $d$, satisfies
\[
 q<\bigl[1+(4H+6)d\bigr]2^d.
\]
\end{proposition}

\begin{proof}
The case $q=0$ is immediate, so assume $q>0$.
Contract each region to one marked vertex.
Also contract every connected component outside the regions, and discard loops.
Call the resulting connected loopless quotient $Q$.
Its $q$ marked vertices have degree at most $d$.
Let $L_v$ be the span of the incident edge-coordinate forms in $\mathcal Z(Q)^*$.
Write $I_v$ for the indicator that this star is zero in a uniform even state of $Q$.

Conditional on the quotient state, the internal cycle-space choices in the contracted regions are independent.
Whenever a region's boundary is zero, its internal even state must be empty for the forest event to hold.
Since that region is cyclic, this has conditional probability at most one half.
Consequently
\[
 \Pr(\xi_G-P\text{ is a forest})
 \le\mathbb E\prod_v(1-I_v/2).
\]
No $L_v$ can be zero, since that would make the right side at most one half.
Also $\dim L_v\le d-1$, by even star parity.

Take a maximal family of these local spaces forming a direct sum.
Let $S$ be its vertex set and $t=|S|$.
The zero-star events for this family are independent, with probabilities at least $2^{1-d}$.
Thus
\[
 \Pr(\xi_G-P\text{ is a forest})
 \le(1-2^{-d})^t.
\]
If $t\ge2^d$, this is at most $e^{-1}<1/2$.
Hence $1\le t<2^d$.
If $S$ contains every marked vertex, the conclusion follows immediately.
Otherwise put $U=\sum_{s\in S}L_s$ and $B=|\delta_Q S|\le dt$.
Maximality gives $L_v\cap U\ne0$ at every remaining marked vertex.

Use the terminal-forest construction from Lemma~\ref{lem:four-terminal-bipartition}.
In each component of $Q-S$, choose a spanning tree and attach a terminal leaf for each edge to $S$.
Prune to paths between terminal leaves, and suppress degree-two vertices.
Except for at most $2B$ marked vertices, all remaining marked vertices lie internally on at most $2B$ resulting tree edges.
Each such vertex is an articulation of $Q-S$ separating its two terminal-bearing sides.
The interior of one resulting tree edge has no branching node and no terminal attachment.
All vertices in that class induce the same bipartition of the attached terminals.

We claim that each such class contains at most $2H+2$ marked vertices.
Suppose a class contains $v_1,\ldots,v_m$ in tree order, with $m\ge2H+3$.
Let $R$ be the component of
\[
 (Q-S)-\{v_1,v_{H+2}\}
\]
containing $v_2,\ldots,v_{H+1}$.
Those vertices belong to the same component because the intervening tree path avoids both deleted vertices.
The articulation property prevents later class vertices from belonging to $R$.

The component $R$ has no neighbor in $S$.
Otherwise it would contain an attached terminal leaf.
That terminal would lie after $v_1$ but before $v_{H+2}$ in their ordered terminal separations.
Their induced terminal bipartitions would then differ, contradicting their membership in the same class.
It follows that every edge of $Q$ leaving $R$ ends at $v_1$ or $v_{H+2}$.
Their combined degree is at most $2d$.

Lift $R$ to $G$ by replacing each contracted vertex by its connected original vertex set.
The lift is connected and contains the $H$ disjoint cyclic regions corresponding to $v_2,\ldots,v_{H+1}$.
Its rank is therefore at least $H$.
On the other side, the later tree segment connects at least $H$ further cyclic regions without entering $R$.
Its lift contains a connected subgraph of rank at least $H$.
Prune leaves from each of these two connected subgraphs, preserving their ranks and obtaining disjoint connected cyclic cores.
The cut around the lift of $R$ separates these cores and has at most $2d$ edges.
Lemma~\ref{lem:random-protected-core-cut} gives a cut of no larger order with connected shores of ranks at least $H$.
This contradicts the hypothesis and proves the claim.

Counting the selected vertices, exceptional vertices, and classes gives
\[
 q\le t+2B+2B(2H+2)
   =t+(4H+6)B
   \le\bigl[1+(4H+6)d\bigr]t.
\]
Since $t<2^d$, the proposition follows.
\end{proof}

For disjoint cycles of length at most $d$ in a cubic graph, the boundary condition holds automatically.
Their total vertex count is therefore less than
\[
 d\bigl[1+(4H+6)d\bigr]2^d.
\]
For each fixed $C$, the fixed-ratio quarter theorem supplies the forbidden-cut hypothesis in an above-half family when $d\le Cg$ and $H=H_{2C}$.
The resulting packing bound is independent of the positive forest excess.
It does not yet supply a common-connected-complement packing or a bound at a girth scale whose ratio to $g$ grows without limit.
